
%
% The document guidelines say the font can be between 10pt and 12pt.
%
\documentclass[10pt]{ClemsonThesis}

\usepackage{amsmath,amsthm,amssymb}
\usepackage{xcolor}
\usepackage{tikz}
\usepackage{float}
\usepackage{braket}
\usepackage{mathtools}
\usepackage{multicol}
\usepackage{enumerate}
\setstretch{2}
%\def\Not{\neg} 
%\def\Not{\overline} 
\newcommand{\Not}[1]{\overline #1}
\DeclareMathOperator{\id}{Id}

\usepackage{comment}
\usepackage{arydshln}

\newcounter{comments}
 \setcounter{comments}{0}
 \newcommand{\fr}[1]{\addtocounter{comments}{1}{\color{blue}[Fabrice \thecomments: #1]}}
 \newcommand{\mm}[1]{\addtocounter{comments}{1}{\color{red}[Matt \thecomments: #1]}}
 
\usepackage{tikz-3dplot}
\usepackage{tikz-network}
\usepackage{mathrsfs}

\usetikzlibrary{arrows}
\usetikzlibrary{arrows.meta}
\usetikzlibrary{decorations.markings}
\usetikzlibrary{decorations.pathreplacing}
\usetikzlibrary{patterns}
\usetikzlibrary{shapes.geometric}
\usetikzlibrary{backgrounds}

% Math theorem environments
\newtheorem{theorem}{Theorem}[section]
\newtheorem{lemma}[theorem]{Lemma}
\newtheorem{proposition}[theorem]{Proposition}
\newtheorem{claim}[theorem]{Claim}
\newtheorem{corollary}[theorem]{Corollary}
\newtheorem{conjecture}[theorem]{Conjecture}

\theoremstyle{definition} 
\newtheorem{definition}[theorem]{Definition}
\newtheorem{example}[theorem]{Example}
\newtheorem{remark}[theorem]{Remark}
\newtheorem{question}[theorem]{Question}

\title{Categories of Boolean networks}
\department{School of Mathematical and Statistical Sciences}
\documentType{Dissertation}
\major{Mathematics}
\degree{Doctor of Philosophy}
\graduationMonth{May}
\graduationYear{2028}
\author{Fabrice Razafimahatratra}
\committeeChair{Dr.\ Matthew Macauley}
\committeeMemberTwo{Dr.\ Keisha Cook}
\committeeMemberThree{Dr.\ Ryann Cartor}
\committeeMemberFour{Dr.\ Cody Stockdale}

\hypersetup{
    colorlinks,
    linkcolor={black},
    citecolor={black},
    filecolor={black},
    urlcolor={black},
    pdftitle={\theTitle},
    pdfauthor={\theAuthor},
    pdfsubject={\theDocumentType},
    pdfkeywords={Clemson University, \theDepartment, \theDocumentType, \theMajor, \theDegree},
    pdfstartpage={1},
}

\usepackage{comment}
\def\<{\langle}
\def\>{\rangle}
\def\longto{\longrightarrow}


%% Sets and Rings
\def\C{\mathbb{C}}
\def\F{\mathbb{F}_2}
\def\N{\mathbb{N}}
\def\Q{\mathbb{Q}}
\def\R{\mathbb{R}}
\def\Z{\mathbb{Z}}

\def\Sum{\sum\limits}
\def\Prod{\prod\limits}
\def\Cap{\mathop\cap\limits}
\def\Cup{\mathop\cup\limits}
\def\Times{\mathop\times\limits}
\def\Oplus{\mathop\oplus\limits}
\def\Lim{\mathop\lim\limits}
\def\del{\partial}

\def\s{\mathbf{s}}
\def\t{\mathbf{t}}

\DeclareMathOperator\Image{Im}
\DeclareMathOperator\Mod{Mod}
\DeclareMathOperator\supp{supp}
\DeclareMathOperator\Sign{sign}
\DeclareMathOperator\XOR{XOR}

%% Categories
\DeclareMathOperator\Hom{Hom}
\DeclareMathOperator\Ob{Ob}
\newcommand{\catname}[1]{{\normalfont\mathbf{#1}}}
\newcommand{\Ab}{\catname{Ab}}
\newcommand{\Field}{\catname{Field}}
\newcommand{\Gph}{\catname{Gph}}
\newcommand{\Grp}{\catname{Grp}}
\newcommand{\LocBool}{\catname{LocBool}}    
\newcommand{\BN}{\catname{BN}}
\newcommand{\FDS}{\catname{FDS}}
\newcommand{\Ring}{\catname{Ring}}
\newcommand{\Rng}{\catname{Rng}}
\newcommand{\Top}{\catname{Top}}
\newcommand{\Boolal}{\catname{BoolAlg}}
\newcommand{\Nilp}{\catname{Nil}}    
\newcommand{\Svl}{\catname{Svl}}
\newcommand{\TruthTable}{\catname{Tbl}}
\newcommand{\fdv}{\catname{FinDimVec}}
\newcommand\cat{\mathcal{C}}
\newcommand\dat{\mathcal{D}}
\newcommand\eat{\mathcal{E}}
\newcommand\fun{\mathcal{F}}
\newcommand{\catset}{\mathbf{Set}}
\newcommand{\Bool}{\mathcal{B}}


\DeclareRobustCommand{\loongrightarrow}{%
  \DOTSB\relbar\joinrel\relbar\joinrel\rightarrow
}

\DeclareRobustCommand{\loongleftarrow}{%
  \DOTSB\leftarrow\joinrel\relbar\joinrel\relbar
}

%% Operators

\def\And{\wedge}
\def\Or{\vee}
%\def\OR{\Or\limits}
%\def\AND{\And\limits}

\newcommand{\join}{\vee}
\newcommand{\meet}{\wedge}
\newcommand{\sloop}{\scalebox{0.8}{\tikz \draw[->] (0,0) arc (530:200:0.15);}}
\DeclareMathOperator\ad{ad}
\def\0{\mathbf{0}}
\def\1{\mathbf{1}}
\def\e{\bm{e}}

\tikzset{%
  Unknown/.tip={Bar[sep]Triangle[angle=45:4pt]},
  %Unknown/.tip={Bar[sep]Triangle[open,angle=45:4pt]},
  tipB/.tip={Bar[sep]Square[open]}
}


\definecolor{m-red}{rgb}{0.8500, 0.3250, 0.0980}
\definecolor{m-blue}{rgb}{0, 0.4470, 0.7410}
\definecolor{m-yellow}{rgb}{.9290, 0.6940, 0.1250}


\begin{document}
%  ============================================================================
    \frontmatter % Begin front matter (pages are numbered with Roman numerals)
%  ============================================================================

    \addtotoc{Title Page}{\maketitle}          % Generate the title page
    \doublespacing                             % Text should be double spaced
    \setcounter{page}{2}                       % Abstract begins on page 2
    \addtotoc{Abstract}{\input{chapters/abstract}}  % Generate the abstract

    % The dedication page is optional.  
    %
    \addtotoc{Dedication}{\input{chapters/dedication.tex}}

    % The acknowledgment page is optional.  
    %
    %\addtotoc{Acknowledgments}{\input{chapters/acknowledgement}}

    \singlespacing                             % Single space the lists
    \tableofcontents \clearpage                % Generate the Table of Contents

    %
    % REMEMBER: Review your caption listings in the genrated lists
    %           and make sure they include '\newline' commands as necessary.
    %           See the README for further information.
    %
   % \addtotoc{List of Tables}{\listoftables}   % Generate the List of Tables
   % \addtotoc{List of Figures}{\listoffigures} % Generate the List of Figures

    %
    % Include other optional lists.  Computer science, for example, would
    % likely include a 'List of Listings' (and would \usepackage{listings}
    % and \renewcommand\lstlistlistingname{List of Listings}).
    %
    %% \addtotoc{List of Listings}{\lstlistoflistings}



%  ===========================================================================
    \mainmatter % Begin main matter (pages are numbered with Arabic numerals)
%  ===========================================================================
    %\doublespacing % Text should be double spaced

    %
    % Here we have each chapter in a separate file.  Name these as you choose,
    % and include them in the order you want them to appear.  Be sure to use
    % the \inputfile command.
    %
    \chapter{Introduction}\label{ch:introduction}
    \inputfile{chapters/introduction.tex}

    \chapter{Maps Between Boolean Networks}\label{ch: maps boolean}
    \inputfile{chapters/Maps_Boolean.tex}

    \chapter{Calculus of Boolean Networks}\label{ch:calculus-boolean}
    \inputfile{chapters/Calculus_Boolean.tex}
   
    \chapter*{Concluding remarks}\label{ch:conclusions}
    \addtotoc{Concluding remarks}{\inputfile{chapters/conclusions.tex}}
    %\inputfile{appendixA.tex}
    %\chapter{Notes}
    %\inputfile{chapters/Notes}
    % The appendices are optional.  This is the format for two or more.
    % If you do not wish to include an appendix, comment out these lines.
    % If you want just one, see the formatting guidelines.
    %
    %\begin{appendices}
       % \begin{subappendices}
            %\inputfile{Appendices/counting_maps.tex}
        %\end{subappendices}
   % \end{appendices}



    \singlespacing                             % Single space the Bibliography


    \bibliographystyle{plain}
    \addtotoc{Bibliography}{\bibliography{References}}
\end{document}
